% Part: lambda-calculus
% Chapter: syntax
% Section: terms

\documentclass[../../../include/open-logic-section]{subfiles}

\begin{document}

\olfileid{lam}{syn}{trm}
\olsection{项}

λ-演算的项是从无穷多个变元 $\Obj{v_0}$、$\Obj{v_1}$、……，符号“$\lambd$”，以及括号出发，归纳地构造出来的。
我们将用 $x$、$y$、$z$、…… 表示变元，用 $M$、$N$、$P$、…… 表示项。

\begin{defn}[项] \ollabel{defn:term}
λ-演算的\emph{项}的集合归纳定义如下：
\begin{enumerate}
  \item \ollabel{defn:term-var} 若 $x$ 是变元，则 $x$ 是项。
  \item \ollabel{defn:term-abs} 若 $x$ 是变元且 $M$ 是项，则 $(\lambd[x][M])$ 是项。
  \item \ollabel{defn:term-app} 若 $M$ 和 $N$ 都是项，则 $(MN)$ 是项。
\end{enumerate}
\end{defn}

如果一个项 $(\lambd[x][M])$ 是依据 \olref{defn:term-abs} 形成的，我们称其为\emph{抽象}的结果，而 $\lambd[x]$ 中的 $x$ 称为\emph{参数}。
依据 \olref{defn:term-app} 形成的项 $(MN)$ 则是\emph{应用}的结果。

上面定义的项是完全括号化的。正如项
$(\lambd[x][((\lambd[x][x])(\lambd[x][(xx)]))])$ 所示，这会变得相当繁琐。
我们将引入省略括号的约定。
不过，原定义使我们很容易判断一个项是如何依据 \olref{defn:term} 构造出来的。
例如，形成项 $(\lambd[x][((\lambd[x][x])(\lambd[x][(xx)]))])$ 的最后一步必定是以~$x$ 为参数的抽象。
它由对项 $((\lambd[x][x])(\lambd[x][(xx)]))$ 进行抽象而来，而该项又是两个项的应用。
这两个项各自又是抽象的结果，依此类推。

\begin{prob}
描述 $(\lambd[g][(\lambd[x][(g (x x))])
  (\lambd[x][(g (x x))])])$ 的构造过程。
\end{prob}

\end{document}